% =============================================================
%  MAT231-F27 Calculus II — Lecture 20
%  Unit 20: Change of variables; polar, cylindrical and spherical
%  coordinates
%
%  Build from the repository root:   npm run materials -- lecture-20
%  This writes lecture-20.pdf (slides, with overlays) and
%  lecture-20-handout.pdf (print version) into this folder.
% =============================================================
\documentclass[10pt]{wosm-slides}

\usepackage{pgfplots}
\pgfplotsset{compat=1.18}

\newcommand{\R}{\mathbb{R}}
\newcommand{\dA}{\, dA}
\newcommand{\dV}{\, dV}
\newcommand{\pd}[2]{\frac{\partial #1}{\partial #2}}

\course{MAT231-F27}{Calculus II}
\title[L20 \textperiodcentered\ Change of variables]{Change of variables}
\subtitle{Polar, cylindrical and spherical coordinates}
\author{Dr.\ Ali Khudiyev}
\institute{WoSM Academy}
\date{Lecture 20\enspace\textperiodcentered\enspace Unit 20\enspace\textperiodcentered\enspace Fall 2027}

\begin{document}

\maketitle

% -------------------------------------------------------------
\begin{frame}{Last time}
  \begin{itemize}
    \item Triple integrals: Riemann sums over boxes; Fubini in any of six orders.
    \item A solid between two surfaces over its shadow; slicing by planes.
    \item Mass, centre of mass, moments of inertia, probabilities as volumes.
  \end{itemize}
  \vskip0.6em
  Discs, balls, cylinders and cones gave square roots in the limits. Today: coordinates in which
  these shapes become boxes, and the factor that every change of coordinates brings with it.
\end{frame}

\begin{frame}{This lecture}
  \tableofcontents
  \vskip1em
  {\small\color{wosm-muted}
  Reading: Strang \& Herman, \emph{Calculus, Volume 3}, Sections 5.3, 5.5 and 5.7.\par}
\end{frame}

% =============================================================
\section{Polar coordinates}

\begin{frame}{Discs are awkward in Cartesian coordinates}
  The volume under $z = 4 - x^2 - y^2$ over the disc $x^2 + y^2 \le 4$ (Lecture 19):
  \[ \int_{-2}^{2} \int_{-\sqrt{4 - x^2}}^{\sqrt{4 - x^2}} (4 - x^2 - y^2)\, dy\, dx . \]
  Possible, but unpleasant.
  \pause
  \vskip0.4em
  In polar coordinates $x = r\cos\theta$, $y = r\sin\theta$ the disc is the rectangle
  \[ 0 \le r \le 2, \qquad 0 \le \theta \le 2\pi , \]
  and the integrand is $4 - r^2$. What is $dA$ in terms of $dr$ and $d\theta$?
\end{frame}

\begin{frame}{The area of a polar rectangle}
  \begin{columns}[c]
    \begin{column}{0.55\textwidth}
      The region $r_0 \le r \le r_0 + \Delta r$, $\theta_0 \le \theta \le \theta_0 + \Delta\theta$ is a piece of
      an annulus. Its area is
      \[ \tfrac12 \bigl((r_0 + \Delta r)^2 - r_0^2\bigr)\, \Delta\theta = \bar r\, \Delta r\, \Delta\theta , \]
      where $\bar r = r_0 + \frac{\Delta r}{2}$ is the middle radius.
      \pause
      It is almost a rectangle with sides $\Delta r$ and $r\,\Delta\theta$. So
      \[ dA = r\, dr\, d\theta . \]
      \pause
      {\small Far from the origin, the same $\Delta\theta$ sweeps out more area: hence the factor $r$.\par}
    \end{column}
    \begin{column}{0.43\textwidth}
      \centering
      \begin{tikzpicture}[scale=1.5]
        \draw[wosm-muted, very thin] (0,0) -- (2.4,0);
        \draw[wosm-rule, dashed] (0,0) -- (25:2.4) (0,0) -- (50:2.4);
        \fill[wosm-accent!25!wosm-bg] (25:1.4) arc (25:50:1.4) -- (50:2.0) arc (50:25:2.0) -- cycle;
        \draw[wosm-text, thick] (25:1.4) arc (25:50:1.4) -- (50:2.0) arc (50:25:2.0) -- cycle;
        \draw[wosm-muted] (0.45,0) arc (0:25:0.45) node[right, font=\scriptsize, xshift=1pt, yshift=-3pt] {$\theta_0$};
        \draw[wosm-brick] (25:0.8) arc (25:50:0.8);
        \node[font=\scriptsize, wosm-brick] at (37.5:0.62) {$\Delta\theta$};
        \node[font=\scriptsize, below right] at (25:1.7) {$\Delta r$};
        \node[font=\scriptsize, above] at (40:2.0) {$r\,\Delta\theta$};
        \fill (0,0) circle (1pt) node[below, font=\scriptsize] {$O$};
      \end{tikzpicture}
    \end{column}
  \end{columns}
\end{frame}

\begin{frame}{Double integrals in polar coordinates}
  \begin{theorem}
    If $D = \{(r\cos\theta, r\sin\theta) : \alpha \le \theta \le \beta, \ h_1(\theta) \le r \le h_2(\theta)\}$
    with $0 \le h_1 \le h_2$ and $\beta - \alpha \le 2\pi$, and $f$ is continuous on $D$, then
    \[ \iint_D f \dA = \int_\alpha^\beta \int_{h_1(\theta)}^{h_2(\theta)} f(r\cos\theta, r\sin\theta)\, r\, dr\, d\theta . \]
  \end{theorem}
  \pause
  \begin{itemize}
    \item Replace $x$ by $r\cos\theta$ and $y$ by $r\sin\theta$ --- and $dA$ by $r\, dr\, d\theta$. Forgetting
      the $r$ is the most common mistake.
    \item The proof: Riemann sums over polar rectangles, whose areas are $\bar r\, \Delta r\, \Delta\theta$.
    \item A general form for any change of coordinates follows shortly.
  \end{itemize}
\end{frame}

\begin{frame}{Examples}
  \begin{example}[The volume from before]
    \[ \int_0^{2\pi} \int_0^2 (4 - r^2)\, r\, dr\, d\theta = 2\pi \Bigl[2r^2 - \frac{r^4}{4}\Bigr]_0^2
      = 2\pi\,(8 - 4) = 8\pi , \]
    as found by slicing in Lecture 19.
  \end{example}
  \pause
  \begin{example}[The annulus $1 \le x^2 + y^2 \le 4$]
    \[ \iint (x^2 + y^2) \dA = \int_0^{2\pi} \int_1^2 r^2 \cdot r\, dr\, d\theta
      = 2\pi \cdot \frac{16 - 1}{4} = \frac{15\pi}{2} . \]
  \end{example}
\end{frame}

\begin{frame}{Example: the cardioid}
  \begin{columns}[c]
    \begin{column}{0.6\textwidth}
      The region inside $r = 1 + \cos\theta$, $0 \le \theta \le 2\pi$:
      \[ \text{area} = \int_0^{2\pi} \!\! \int_0^{1 + \cos\theta} \! r\, dr\, d\theta
         = \frac12 \int_0^{2\pi} (1 + \cos\theta)^2\, d\theta . \]
      \pause
      Expand: $\int_0^{2\pi} (1 + 2\cos\theta + \cos^2\theta)\, d\theta = 2\pi + 0 + \pi$. \ So the area is
      $\frac{3\pi}{2}$.
      \pause
      \vskip0.4em
      {\small The formula $\frac12 \int r(\theta)^2\, d\theta$ for the area of a region bounded by a polar
      curve is the special case $f = 1$, $h_1 = 0$.\par}
    \end{column}
    \begin{column}{0.38\textwidth}
      \centering
      \begin{tikzpicture}[scale=1.25]
        \draw[wosm-muted, very thin] (-0.6,0) -- (2.3,0) (0,-1.5) -- (0,1.5);
        \fill[wosm-accent!20!wosm-bg, domain=0:360, samples=120, variable=\t] plot ({(1 + cos(\t))*cos(\t)}, {(1 + cos(\t))*sin(\t)}) -- cycle;
        \draw[wosm-accent, thick, domain=0:360, samples=120, variable=\t] plot ({(1 + cos(\t))*cos(\t)}, {(1 + cos(\t))*sin(\t)});
        \draw[wosm-brick, thick] (0,0) -- ({(1 + cos(50))*cos(50)}, {(1 + cos(50))*sin(50)});
        \node[font=\scriptsize, wosm-brick, above left] at ({0.5*(1 + cos(50))*cos(50)}, {0.5*(1 + cos(50))*sin(50)}) {$r$};
      \end{tikzpicture}
      \par\small $r = 1 + \cos\theta$
    \end{column}
  \end{columns}
\end{frame}

% =============================================================
\section{The change of variables formula}

\begin{frame}{Substitution in one variable}
  For a $C^1$ function $g$ that is increasing or decreasing on $[a, b]$:
  \[ \int_{g([a, b])} f(x)\, dx = \int_a^b f\bigl(g(u)\bigr)\, |g'(u)|\, du \qquad \text{(Calculus I, Unit 21).} \]
  \pause
  \begin{itemize}
    \item $|g'(u)|$ is the \textbf{stretch factor}: $g$ maps a short interval of length $\Delta u$ to one of
      length about $|g'(u)|\, \Delta u$.
    \item In two variables a map $T(u, v) = \bigl(x(u, v), y(u, v)\bigr)$ sends small rectangles to small
      (curved) parallelograms. The factor is the ratio of their areas.
  \end{itemize}
\end{frame}

\begin{frame}{The Jacobian as an area factor}
  $T$ maps the small rectangle $[u, u + \Delta u] \times [v, v + \Delta v]$ to a region that is almost the
  parallelogram spanned by $T_u\, \Delta u$ and $T_v\, \Delta v$ \ (since $T(u + \Delta u, v) - T(u, v) \approx T_u\, \Delta u$).
  \begin{columns}[c]
    \begin{column}{0.52\textwidth}
      \pause
      Its area is $|\det(T_u, T_v)|\, \Delta u\, \Delta v$ (Lecture 12). With the \textbf{Jacobian}
      \[ \pd{(x, y)}{(u, v)} = \det DT = \det\begin{pmatrix} x_u & x_v \\ y_u & y_v \end{pmatrix} , \]
      \[ dA = \Bigl| \pd{(x, y)}{(u, v)} \Bigr|\, du\, dv . \]
    \end{column}
    \begin{column}{0.46\textwidth}
      \centering
      \begin{tikzpicture}[scale=1.3]
        \draw[wosm-muted, -latex] (0,0) -- (1.7,0) node[right, font=\scriptsize] {$u$};
        \draw[wosm-muted, -latex] (0,0) -- (0,1.7) node[above, font=\scriptsize] {$v$};
        \foreach \s in {0.3, 0.7, 1.1, 1.5} {
          \draw[wosm-rule] (\s,0.3) -- (\s,1.5);
          \draw[wosm-rule] (0.3,\s) -- (1.5,\s);
        }
        \fill[wosm-accent!35!wosm-bg] (0.7,0.7) rectangle (1.1,1.1);
        \draw[wosm-text] (0.7,0.7) rectangle (1.1,1.1);
        \draw[wosm-brick, -latex] (1.75,0.9) to[bend left=20] node[above, font=\scriptsize] {$T$} (2.35,0.9);
        \begin{scope}[shift={(2.3,-0.2)}]
          \foreach \s in {0.3, 0.7, 1.1, 1.5} {
            \draw[wosm-rule, domain=0.3:1.5, samples=20, variable=\t] plot ({\s + 0.25*\t*\t}, {\t + 0.2*\s*\s});
            \draw[wosm-rule, domain=0.3:1.5, samples=20, variable=\t] plot ({\t + 0.25*\s*\s}, {\s + 0.2*\t*\t});
          }
          \fill[wosm-accent!35!wosm-bg] plot[domain=0.7:1.1, samples=10, variable=\t] ({\t + 0.25*0.49}, {0.7 + 0.2*\t*\t})
            -- plot[domain=0.7:1.1, samples=10, variable=\t] ({1.1 + 0.25*\t*\t}, {\t + 0.2*1.21})
            -- plot[domain=1.1:0.7, samples=10, variable=\t] ({\t + 0.25*1.21}, {1.1 + 0.2*\t*\t})
            -- plot[domain=1.1:0.7, samples=10, variable=\t] ({0.7 + 0.25*\t*\t}, {\t + 0.2*0.49}) -- cycle;
          \draw[wosm-brick, thick, -latex] (0.8225,0.798) -- (1.2225,0.910) node[below right, font=\scriptsize] {$T_u \Delta u$};
          \draw[wosm-brick, thick, -latex] (0.8225,0.798) -- (0.9625,1.198) node[above, font=\scriptsize] {$T_v \Delta v$};
        \end{scope}
      \end{tikzpicture}
    \end{column}
  \end{columns}
\end{frame}

\begin{frame}{The change of variables theorem}
  \begin{theorem}
    Let $T\colon S \to D$ be a $C^1$ map from a region $S$ in the $uv$-plane onto a region $D$ in the
    $xy$-plane, one-to-one with $\det DT \ne 0$ except possibly on the boundary of $S$. Then for
    continuous $f$
    \[ \iint_D f(x, y) \dA = \iint_S f\bigl(T(u, v)\bigr)\, \Bigl| \pd{(x, y)}{(u, v)} \Bigr|\, du\, dv . \]
  \end{theorem}
  \pause
  \begin{itemize}
    \item Polar: \ $\pd{(x, y)}{(r, \theta)} = r$ (Lecture 17), which is $0$ only on the boundary $r = 0$.
    \item The same holds in $\R^n$, with $|\det DT|$ the volume factor.
    \item Choose $T$ so that $S$ is simple (a rectangle), or the integrand is.
  \end{itemize}
\end{frame}

\begin{frame}{Example: a linear change}
  $D$ is the square bounded by the lines $x + y = 0$, $x + y = 2$, $x - y = 0$, $x - y = 2$.
  Compute $\iint_D (x^2 - y^2) \dA$.
  \pause
  \begin{itemize}
    \item Put $u = x + y$, $v = x - y$. Then $S = [0, 2] \times [0, 2]$, and $x^2 - y^2 = uv$.
    \item $x = \frac{u + v}{2}$, $y = \frac{u - v}{2}$: \
      $\pd{(x, y)}{(u, v)} = \det\begin{pmatrix} 1/2 & 1/2 \\ 1/2 & -1/2 \end{pmatrix} = -\frac12$.
  \end{itemize}
  \pause
  \[ \iint_D (x^2 - y^2) \dA = \int_0^2 \int_0^2 uv \cdot \frac12\, du\, dv = \frac12 \cdot 2 \cdot 2 = 2 . \]
  \pause
  \begin{example}[The ellipse $\frac{x^2}{a^2} + \frac{y^2}{b^2} \le 1$]
    $x = au$, $y = bv$ maps the unit disc onto it, with Jacobian $ab$. So its area is $ab \cdot \pi$.
  \end{example}
\end{frame}

\begin{frame}{Using the inverse Jacobian}
  Often $u$, $v$ are given as functions of $x$, $y$. By the inverse function theorem (Lecture 17),
  $\pd{(x, y)}{(u, v)} = 1 \big/ \pd{(u, v)}{(x, y)}$.
  \pause
  \vskip0.5em
  \textbf{Example.} The region $D$ in the first quadrant between $xy = 1$, $xy = 3$, $y = x$, $y = 2x$.
  \begin{itemize}
    \item Put $u = xy \in [1, 3]$ and $v = y/x \in [1, 2]$: \ $D$ becomes a rectangle.
    \pause
    \item $\pd{(u, v)}{(x, y)} = \det\begin{pmatrix} y & x \\ -y/x^2 & 1/x \end{pmatrix} = \frac{2y}{x} = 2v$, \
      so $dA = \dfrac{du\, dv}{2v}$.
    \pause
    \item $\displaystyle \operatorname{area}(D) = \int_1^3 \int_1^2 \frac{dv}{2v}\, du = 2 \cdot \frac{\ln 2}{2} = \ln 2$.
  \end{itemize}
  \pause
  {\small No need to solve for $x$ and $y$ (which gives $x = \sqrt{u/v}$, $y = \sqrt{uv}$).\par}
\end{frame}

% =============================================================
\section{Cylindrical coordinates}

\begin{frame}{Cylindrical coordinates}
  Polar coordinates in the $xy$-plane, and $z$ unchanged:
  \[ x = r\cos\theta, \quad y = r\sin\theta, \quad z = z; \qquad
     \pd{(x, y, z)}{(r, \theta, z)} = r, \quad dV = r\, dr\, d\theta\, dz . \]
  Good for solids with an axis of symmetry: cylinders, cones, paraboloids.
  \pause
  \begin{example}[The paraboloid solid, a third time]
    $r^2 \le z \le 4$ over $0 \le r \le 2$:
    \[ \int_0^{2\pi} \int_0^2 \int_{r^2}^4 r\, dz\, dr\, d\theta = 2\pi \int_0^2 (4r - r^3)\, dr = 8\pi . \]
  \end{example}
\end{frame}

\begin{frame}{Examples in cylindrical coordinates}
  \begin{example}[The cone $\sqrt{x^2 + y^2} \le z \le 1$, that is $r \le z \le 1$]
    \[ \int_0^{2\pi} \int_0^1 \int_r^1 r\, dz\, dr\, d\theta = 2\pi \int_0^1 (r - r^2)\, dr = \frac\pi3
       = \frac13 \cdot \text{base} \cdot \text{height} . \]
  \end{example}
  \pause
  \begin{example}[A solid cylinder: radius $a$, height $h$, density $1$]
    Its moment of inertia about its axis ($M = \pi a^2 h$):
    \[ I = \int_0^{2\pi} \int_0^a \int_0^h r^2 \cdot r\, dz\, dr\, d\theta = \frac{\pi a^4 h}{2} = \frac12 M a^2 . \]
  \end{example}
\end{frame}

% =============================================================
\section{Spherical coordinates}

\begin{frame}{Spherical coordinates}
  \[ x = \rho\sin\varphi\cos\theta, \qquad y = \rho\sin\varphi\sin\theta, \qquad z = \rho\cos\varphi \]
  \begin{columns}[c]
    \begin{column}{0.6\textwidth}
      \begin{itemize}
        \item $\rho \ge 0$: the distance to $O$;
        \item $0 \le \varphi \le \pi$: the angle from the $z$-axis;
        \item $0 \le \theta \le 2\pi$: as in polar coordinates.
      \end{itemize}
      \pause
      The Jacobian determinant is $\rho^2 \sin\varphi \ge 0$ (a small spherical box has edges
      $\Delta\rho$, $\rho\,\Delta\varphi$, $\rho\sin\varphi\,\Delta\theta$):
      \[ dV = \rho^2 \sin\varphi\, d\rho\, d\varphi\, d\theta . \]
    \end{column}
    \begin{column}{0.38\textwidth}
      \centering
      \begin{tikzpicture}[x={(-0.6cm,-0.42cm)}, y={(1cm,0cm)}, z={(0cm,1cm)}, scale=1.3]
        \draw[wosm-muted, -latex] (0,0,0) -- (1.9,0,0) node[below left, font=\scriptsize] {$x$};
        \draw[wosm-muted, -latex] (0,0,0) -- (0,1.9,0) node[right, font=\scriptsize] {$y$};
        \draw[wosm-muted, -latex] (0,0,0) -- (0,0,2.1) node[above, font=\scriptsize] {$z$};
        \draw[wosm-rule, dashed] (0,0,0) -- (0.643,1.113,0) -- (0.643,1.113,1.532);
        \draw[wosm-brick, thick] (0,0,0) -- (0.643,1.113,1.532) node[midway, below right, font=\scriptsize] {$\rho$};
        \fill[wosm-brick] (0.643,1.113,1.532) circle (1.3pt) node[right, font=\scriptsize] {$P$};
        \draw[wosm-accent, domain=0:40, samples=12, variable=\t] plot ({0.6*sin(\t)*cos(60)}, {0.6*sin(\t)*sin(60)}, {0.6*cos(\t)});
        \node[font=\scriptsize, wosm-accent] at (0,0.12,0.88) {$\varphi$};
        \draw[wosm-accent, domain=0:60, samples=12, variable=\t] plot ({0.55*cos(\t)}, {0.55*sin(\t)}, 0);
        \node[font=\scriptsize, wosm-accent] at (0.8,0.5,0) {$\theta$};
      \end{tikzpicture}
    \end{column}
  \end{columns}
\end{frame}

\begin{frame}{Examples in spherical coordinates}
  \begin{example}[The ball of radius $R$]
    \[ \int_0^{2\pi} \int_0^\pi \int_0^R \rho^2 \sin\varphi\, d\rho\, d\varphi\, d\theta
      = 2\pi \cdot 2 \cdot \frac{R^3}{3} = \frac43 \pi R^3 . \]
  \end{example}
  \pause
  \begin{example}[$\iiint e^{(x^2 + y^2 + z^2)^{3/2}} \dV$ over the unit ball]
    \[ 2\pi \cdot 2 \cdot \int_0^1 e^{\rho^3} \rho^2\, d\rho = 4\pi \cdot \frac{e - 1}{3} \]
    --- hopeless in Cartesian coordinates.
  \end{example}
\end{frame}

\begin{frame}{Example: the ice-cream cone}
  \begin{columns}[c]
    \begin{column}{0.6\textwidth}
      The solid inside the sphere $\rho \le 1$ and within the cone $\varphi \le \frac\pi4$, that is
      $z \ge \sqrt{x^2 + y^2}$.
      \pause
      \[ \int_0^{2\pi} \int_0^{\pi/4} \int_0^1 \rho^2 \sin\varphi\, d\rho\, d\varphi\, d\theta \]
      \[ = 2\pi \Bigl(1 - \frac{\sqrt2}{2}\Bigr) \frac13 = \frac{\pi\,(2 - \sqrt2)}{3} \approx 0.613 . \]
      \pause
      {\small In spherical coordinates the region is a box: $[0, 1] \times [0, \frac\pi4] \times [0, 2\pi]$.\par}
    \end{column}
    \begin{column}{0.38\textwidth}
      \centering
      \begin{tikzpicture}[scale=1.5]
        \draw[wosm-muted, very thin] (-1.2,0) -- (1.2,0) (0,-0.2) -- (0,1.25);
        \fill[wosm-accent!25!wosm-bg] (0,0) -- (45:1) arc (45:135:1) -- cycle;
        \draw[wosm-accent, thick] (0,0) -- (45:1) arc (45:135:1) -- cycle;
        \draw[wosm-rule, dashed] (1,0) arc (0:45:1) (135:1) arc (135:180:1);
        \draw[wosm-brick] (0,0.35) arc (90:45:0.35);
        \node[font=\scriptsize, wosm-brick] at (67:0.5) {$\frac\pi4$};
      \end{tikzpicture}
      \par\small A cross-section through the $z$-axis.
    \end{column}
  \end{columns}
\end{frame}

\begin{frame}{Example: the centroid of a hemisphere}
  The solid hemisphere $\rho \le R$, $0 \le \varphi \le \frac\pi2$ has volume $\frac23 \pi R^3$, and
  $\bar x = \bar y = 0$ by symmetry.
  \pause
  \[ \iiint z \dV = \int_0^{2\pi} \int_0^{\pi/2} \int_0^R (\rho\cos\varphi)\, \rho^2 \sin\varphi\,
     d\rho\, d\varphi\, d\theta = 2\pi \cdot \frac12 \cdot \frac{R^4}{4} = \frac{\pi R^4}{4} . \]
  \pause
  \[ \bar z = \frac{\pi R^4 / 4}{2\pi R^3 / 3} = \frac{3R}{8} . \]
  \pause
  {\small The centroid is below the middle height $\frac R2$: most of the volume is near the flat base.\par}
\end{frame}

\begin{frame}{Which coordinates?}
  \begin{center}
    \begin{tabular}{@{}ll@{}}
      \toprule
      Region & Coordinates \\
      \midrule
      discs, annuli, sectors, polar curves & polar \\
      cylinders, cones and paraboloids around an axis & cylindrical \\
      balls, spherical shells, cones with vertex at $O$ & spherical \\
      parallelograms, regions between level curves & a linear or adapted change \\
      \bottomrule
    \end{tabular}
  \end{center}
  \pause
  \begin{itemize}
    \item Look at the region first, then at the integrand: $x^2 + y^2$ suggests polar or cylindrical,
      $x^2 + y^2 + z^2$ spherical.
    \item Always include the Jacobian factor: $r$, $r$, $\rho^2\sin\varphi$, or $|\det DT|$.
  \end{itemize}
\end{frame}

% =============================================================
\section{Summary and practice}

\begin{frame}{Summary}
  \begin{itemize}
    \item \textbf{Change of variables}: $\iint_D f \dA = \iint_S f(T(u, v))\,
      \bigl|\frac{\partial(x, y)}{\partial(u, v)}\bigr|\, du\, dv$; the Jacobian determinant is the area
      (volume) factor.
    \item Polar: $dA = r\, dr\, d\theta$. \ Cylindrical: $dV = r\, dr\, d\theta\, dz$. \
      Spherical: $dV = \rho^2 \sin\varphi\, d\rho\, d\varphi\, d\theta$.
    \item Choose coordinates in which the region is a rectangle or box, or the integrand simplifies.
    \item If $u, v$ are given in terms of $x, y$, use $\frac{\partial(x, y)}{\partial(u, v)} = 1 \big/ \frac{\partial(u, v)}{\partial(x, y)}$.
  \end{itemize}
\end{frame}

\begin{frame}{Practice}
  \small
  \begin{enumerate}
    \item Compute $\iint \sqrt{x^2 + y^2} \dA$ over the disc $x^2 + y^2 \le 9$.
    \item Find the area inside the circle $r = 2\sin\theta$.
    \item Compute $\iint_D (x + y)^2 \dA$, where $D$ is the square $|x| + |y| \le 1$.
    \item Find the volume of the ellipsoid $\frac{x^2}{a^2} + \frac{y^2}{b^2} + \frac{z^2}{c^2} \le 1$.
    \item Find the volume inside the cylinder $x^2 + y^2 \le 1$, above $z = 0$ and below
      $z = 4 - x^2 - y^2$.
    \item Find the mass of the unit ball if the density is the distance to the centre.
    \item Find the volume between the spheres of radius $1$ and $2$ (centred at $O$) within the cone
      $\varphi \le \frac\pi3$.
  \end{enumerate}
\end{frame}

\begin{frame}{Answers}
  \small
  \begin{enumerate}
    \item $\int_0^{2\pi} \int_0^3 r \cdot r\, dr\, d\theta = 2\pi \cdot 9 = 18\pi$.
    \item $\frac12 \int_0^\pi 4\sin^2\theta\, d\theta = \pi$ \ (a circle of radius $1$).
    \item $u = x + y$, $v = x - y$ range over $[-1, 1]^2$, with $|J| = \frac12$:
      $\frac12 \int_{-1}^1 \int_{-1}^1 u^2\, du\, dv = \frac23$.
    \item $x = au$, $y = bv$, $z = cw$ maps the unit ball onto it, with Jacobian $abc$: \ $\frac43 \pi abc$.
    \item $\int_0^{2\pi} \int_0^1 (4 - r^2)\, r\, dr\, d\theta = 2\pi\bigl(2 - \frac14\bigr) = \frac{7\pi}{2}$.
    \item $\int_0^{2\pi} \int_0^\pi \int_0^1 \rho \cdot \rho^2 \sin\varphi\, d\rho\, d\varphi\, d\theta = 4\pi \cdot \frac14 = \pi$.
    \item $2\pi \bigl(1 - \cos\frac\pi3\bigr) \cdot \frac{2^3 - 1}{3} = \frac{7\pi}{3}$.
  \end{enumerate}
\end{frame}

\begin{frame}{Next time}
  \textbf{Lecture 21 --- The Gaussian integral; improper multiple integrals} (Unit 21)
  \begin{itemize}
    \item $\int_{-\infty}^\infty e^{-x^2}\, dx = \sqrt\pi$, by a double integral in polar coordinates
    \item When do integrals over unbounded regions, or of unbounded functions, converge?
    \item The normal distribution, and the volume of a ball in $n$ dimensions
  \end{itemize}
  \vskip1em
  {\small
  Reading: Marsden \& Tromba, \emph{Vector Calculus}, Section 6.4.\\[0.4em]
  Questions: the module channel on Discord, or office hours,
  Monday--Friday, 10:00--17:00 (Baku time, UTC+4),
  \href{mailto:ali.khudiyev@wosm.academy}{ali.khudiyev@wosm.academy}.\par}
\end{frame}

\end{document}
